% =====================================================================
%  Introduction : importée par main.tex avec % =====================================================================
%  Introduction : importée par main.tex avec % =====================================================================
%  Introduction : importée par main.tex avec % =====================================================================
%  Introduction : importée par main.tex avec \include{introduction}.
%  Section non numérotée, ajoutée au sommaire. Pas de figure flottante ici
%  (la numérotation des figures suit celle des chapitres) : le plan est un
%  dessin TikZ en place.
% =====================================================================

% Macros de Dirac (déjà définies dans main.tex ; ce bloc évite l'erreur
% « Undefined control sequence \ket » si main.tex est une ancienne version).
\providecommand{\ket}[1]{|#1\rangle}
\providecommand{\bra}[1]{\langle #1|}
\providecommand{\braket}[2]{\langle #1|#2\rangle}

\phantomsection
\section*{Introduction}
\addcontentsline{toc}{section}{Introduction}

Ce document rassemble les notes du cours d'informatique quantique de \docprof{}
(4A FISE SF). Il part de la mécanique ondulatoire (fonction d'onde de Schrödinger,
interprétation de Born) et en tire, pas à pas, les outils de l'informatique
quantique : états d'un \gls{qubit}, \glslink{hilbert}{espace de Hilbert}, opérateurs,
matrice de densité, portes quantiques, systèmes de plusieurs qubits,
\gls{intrication} et téléportation. Tous les calculs sont détaillés.

\begin{center}
\begin{tikzpicture}[>=Stealth, font=\scriptsize,
    boite/.style={draw=insarouge, fill=insarouge!6, thick, rounded corners=2pt,
                  text width=2.02cm, minimum height=1.55cm, align=center},
    num/.style={circle, fill=insarouge, text=white, font=\small\bfseries, minimum size=5.5mm, inner sep=0pt}]
  \foreach \i/\txt/\idee in {
      1/{Histoire et mécanique quantique}/{onde de matière},
      2/{Notation de Dirac et espace de Hilbert}/{états = vecteurs},
      3/{Opérateurs, matrice de densité, sphère de Bloch, produit tensoriel}/{états mixtes, Bloch},
      4/{Portes quantiques et systèmes multi-qubits}/{portes unitaires},
      5/{États quantiques, intrication, états de Bell}/{corrélations, téléportation},
      6/{Exemples et exercices corrigés}/{s'entraîner},
      7/{QCM sur le cours}/{s'évaluer}} {
    \node[boite] (b\i) at ({(\i-1)*2.36},0) {\txt};
    \node[num] at ({(\i-1)*2.36},1.2) {\i};
    \node[text=insagris, font=\scriptsize\itshape, text width=2.1cm, align=center] at ({(\i-1)*2.36},-1.25) {\idee};
  }
  \foreach \i/\j in {1/2, 2/3, 3/4, 4/5, 5/6, 6/7}
    \draw[->, insarouge, thick] (b\i.east) -- (b\j.west);
\end{tikzpicture}
\end{center}

Chaque chapitre se termine par un \emph{récapitulatif} : schéma de synthèse, points à retenir,
tableau de formules renvoyant aux sections, et liste de savoir-faire. Le chapitre~\ref{chap:exercices}
rassemble des exemples et exercices corrigés, et le chapitre~\ref{chap:qcm} un QCM de 100 questions
avec corrigé détaillé. Les termes techniques sont définis dans
le glossaire, en fin de document ; les ouvrages de référence figurent dans la bibliographie
\parencite[par exemple][]{nielsen2010,cohen1973}.

\paragraph{Prérequis et conventions.} Algèbre linéaire (matrices, valeurs propres),
nombres complexes, notions d'ondes. Le symbole $\ii$ désigne l'unité imaginaire, $\ket{0}$
et $\ket{1}$ la base de calcul d'un qubit, et $\hbar=h/2\pi$.
.
%  Section non numérotée, ajoutée au sommaire. Pas de figure flottante ici
%  (la numérotation des figures suit celle des chapitres) : le plan est un
%  dessin TikZ en place.
% =====================================================================

% Macros de Dirac (déjà définies dans main.tex ; ce bloc évite l'erreur
% « Undefined control sequence \ket » si main.tex est une ancienne version).
\providecommand{\ket}[1]{|#1\rangle}
\providecommand{\bra}[1]{\langle #1|}
\providecommand{\braket}[2]{\langle #1|#2\rangle}

\phantomsection
\section*{Introduction}
\addcontentsline{toc}{section}{Introduction}

Ce document rassemble les notes du cours d'informatique quantique de \docprof{}
(4A FISE SF). Il part de la mécanique ondulatoire (fonction d'onde de Schrödinger,
interprétation de Born) et en tire, pas à pas, les outils de l'informatique
quantique : états d'un \gls{qubit}, \glslink{hilbert}{espace de Hilbert}, opérateurs,
matrice de densité, portes quantiques, systèmes de plusieurs qubits,
\gls{intrication} et téléportation. Tous les calculs sont détaillés.

\begin{center}
\begin{tikzpicture}[>=Stealth, font=\scriptsize,
    boite/.style={draw=insarouge, fill=insarouge!6, thick, rounded corners=2pt,
                  text width=2.02cm, minimum height=1.55cm, align=center},
    num/.style={circle, fill=insarouge, text=white, font=\small\bfseries, minimum size=5.5mm, inner sep=0pt}]
  \foreach \i/\txt/\idee in {
      1/{Histoire et mécanique quantique}/{onde de matière},
      2/{Notation de Dirac et espace de Hilbert}/{états = vecteurs},
      3/{Opérateurs, matrice de densité, sphère de Bloch, produit tensoriel}/{états mixtes, Bloch},
      4/{Portes quantiques et systèmes multi-qubits}/{portes unitaires},
      5/{États quantiques, intrication, états de Bell}/{corrélations, téléportation},
      6/{Exemples et exercices corrigés}/{s'entraîner},
      7/{QCM sur le cours}/{s'évaluer}} {
    \node[boite] (b\i) at ({(\i-1)*2.36},0) {\txt};
    \node[num] at ({(\i-1)*2.36},1.2) {\i};
    \node[text=insagris, font=\scriptsize\itshape, text width=2.1cm, align=center] at ({(\i-1)*2.36},-1.25) {\idee};
  }
  \foreach \i/\j in {1/2, 2/3, 3/4, 4/5, 5/6, 6/7}
    \draw[->, insarouge, thick] (b\i.east) -- (b\j.west);
\end{tikzpicture}
\end{center}

Chaque chapitre se termine par un \emph{récapitulatif} : schéma de synthèse, points à retenir,
tableau de formules renvoyant aux sections, et liste de savoir-faire. Le chapitre~\ref{chap:exercices}
rassemble des exemples et exercices corrigés, et le chapitre~\ref{chap:qcm} un QCM de 100 questions
avec corrigé détaillé. Les termes techniques sont définis dans
le glossaire, en fin de document ; les ouvrages de référence figurent dans la bibliographie
\parencite[par exemple][]{nielsen2010,cohen1973}.

\paragraph{Prérequis et conventions.} Algèbre linéaire (matrices, valeurs propres),
nombres complexes, notions d'ondes. Le symbole $\ii$ désigne l'unité imaginaire, $\ket{0}$
et $\ket{1}$ la base de calcul d'un qubit, et $\hbar=h/2\pi$.
.
%  Section non numérotée, ajoutée au sommaire. Pas de figure flottante ici
%  (la numérotation des figures suit celle des chapitres) : le plan est un
%  dessin TikZ en place.
% =====================================================================

% Macros de Dirac (déjà définies dans main.tex ; ce bloc évite l'erreur
% « Undefined control sequence \ket » si main.tex est une ancienne version).
\providecommand{\ket}[1]{|#1\rangle}
\providecommand{\bra}[1]{\langle #1|}
\providecommand{\braket}[2]{\langle #1|#2\rangle}

\phantomsection
\section*{Introduction}
\addcontentsline{toc}{section}{Introduction}

Ce document rassemble les notes du cours d'informatique quantique de \docprof{}
(4A FISE SF). Il part de la mécanique ondulatoire (fonction d'onde de Schrödinger,
interprétation de Born) et en tire, pas à pas, les outils de l'informatique
quantique : états d'un \gls{qubit}, \glslink{hilbert}{espace de Hilbert}, opérateurs,
matrice de densité, portes quantiques, systèmes de plusieurs qubits,
\gls{intrication} et téléportation. Tous les calculs sont détaillés.

\begin{center}
\begin{tikzpicture}[>=Stealth, font=\scriptsize,
    boite/.style={draw=insarouge, fill=insarouge!6, thick, rounded corners=2pt,
                  text width=2.02cm, minimum height=1.55cm, align=center},
    num/.style={circle, fill=insarouge, text=white, font=\small\bfseries, minimum size=5.5mm, inner sep=0pt}]
  \foreach \i/\txt/\idee in {
      1/{Histoire et mécanique quantique}/{onde de matière},
      2/{Notation de Dirac et espace de Hilbert}/{états = vecteurs},
      3/{Opérateurs, matrice de densité, sphère de Bloch, produit tensoriel}/{états mixtes, Bloch},
      4/{Portes quantiques et systèmes multi-qubits}/{portes unitaires},
      5/{États quantiques, intrication, états de Bell}/{corrélations, téléportation},
      6/{Exemples et exercices corrigés}/{s'entraîner},
      7/{QCM sur le cours}/{s'évaluer}} {
    \node[boite] (b\i) at ({(\i-1)*2.36},0) {\txt};
    \node[num] at ({(\i-1)*2.36},1.2) {\i};
    \node[text=insagris, font=\scriptsize\itshape, text width=2.1cm, align=center] at ({(\i-1)*2.36},-1.25) {\idee};
  }
  \foreach \i/\j in {1/2, 2/3, 3/4, 4/5, 5/6, 6/7}
    \draw[->, insarouge, thick] (b\i.east) -- (b\j.west);
\end{tikzpicture}
\end{center}

Chaque chapitre se termine par un \emph{récapitulatif} : schéma de synthèse, points à retenir,
tableau de formules renvoyant aux sections, et liste de savoir-faire. Le chapitre~\ref{chap:exercices}
rassemble des exemples et exercices corrigés, et le chapitre~\ref{chap:qcm} un QCM de 100 questions
avec corrigé détaillé. Les termes techniques sont définis dans
le glossaire, en fin de document ; les ouvrages de référence figurent dans la bibliographie
\parencite[par exemple][]{nielsen2010,cohen1973}.

\paragraph{Prérequis et conventions.} Algèbre linéaire (matrices, valeurs propres),
nombres complexes, notions d'ondes. Le symbole $\ii$ désigne l'unité imaginaire, $\ket{0}$
et $\ket{1}$ la base de calcul d'un qubit, et $\hbar=h/2\pi$.
.
%  Section non numérotée, ajoutée au sommaire. Pas de figure flottante ici
%  (la numérotation des figures suit celle des chapitres) : le plan est un
%  dessin TikZ en place.
% =====================================================================

% Macros de Dirac (déjà définies dans main.tex ; ce bloc évite l'erreur
% « Undefined control sequence \ket » si main.tex est une ancienne version).
\providecommand{\ket}[1]{|#1\rangle}
\providecommand{\bra}[1]{\langle #1|}
\providecommand{\braket}[2]{\langle #1|#2\rangle}

\phantomsection
\section*{Introduction}
\addcontentsline{toc}{section}{Introduction}

Ce document rassemble les notes du cours d'informatique quantique de \docprof{}
(4A FISE SF). Il part de la mécanique ondulatoire (fonction d'onde de Schrödinger,
interprétation de Born) et en tire, pas à pas, les outils de l'informatique
quantique : états d'un \gls{qubit}, \glslink{hilbert}{espace de Hilbert}, opérateurs,
matrice de densité, portes quantiques, systèmes de plusieurs qubits,
\gls{intrication} et téléportation. Tous les calculs sont détaillés.

\begin{center}
\begin{tikzpicture}[>=Stealth, font=\scriptsize,
    boite/.style={draw=insarouge, fill=insarouge!6, thick, rounded corners=2pt,
                  text width=2.02cm, minimum height=1.55cm, align=center},
    num/.style={circle, fill=insarouge, text=white, font=\small\bfseries, minimum size=5.5mm, inner sep=0pt}]
  \foreach \i/\txt/\idee in {
      1/{Histoire et mécanique quantique}/{onde de matière},
      2/{Notation de Dirac et espace de Hilbert}/{états = vecteurs},
      3/{Opérateurs, matrice de densité, sphère de Bloch, produit tensoriel}/{états mixtes, Bloch},
      4/{Portes quantiques et systèmes multi-qubits}/{portes unitaires},
      5/{États quantiques, intrication, états de Bell}/{corrélations, téléportation},
      6/{Exemples et exercices corrigés}/{s'entraîner},
      7/{QCM sur le cours}/{s'évaluer}} {
    \node[boite] (b\i) at ({(\i-1)*2.36},0) {\txt};
    \node[num] at ({(\i-1)*2.36},1.2) {\i};
    \node[text=insagris, font=\scriptsize\itshape, text width=2.1cm, align=center] at ({(\i-1)*2.36},-1.25) {\idee};
  }
  \foreach \i/\j in {1/2, 2/3, 3/4, 4/5, 5/6, 6/7}
    \draw[->, insarouge, thick] (b\i.east) -- (b\j.west);
\end{tikzpicture}
\end{center}

Chaque chapitre se termine par un \emph{récapitulatif} : schéma de synthèse, points à retenir,
tableau de formules renvoyant aux sections, et liste de savoir-faire. Le chapitre~\ref{chap:exercices}
rassemble des exemples et exercices corrigés, et le chapitre~\ref{chap:qcm} un QCM de 100 questions
avec corrigé détaillé. Les termes techniques sont définis dans
le glossaire, en fin de document ; les ouvrages de référence figurent dans la bibliographie
\parencite[par exemple][]{nielsen2010,cohen1973}.

\paragraph{Prérequis et conventions.} Algèbre linéaire (matrices, valeurs propres),
nombres complexes, notions d'ondes. Le symbole $\ii$ désigne l'unité imaginaire, $\ket{0}$
et $\ket{1}$ la base de calcul d'un qubit, et $\hbar=h/2\pi$.
